\documentclass[12pt,a4paper]{article}
\usepackage[utf8]{inputenc}
\usepackage[T1]{fontenc}
\usepackage{amsmath,amssymb,amsthm}
\usepackage{mathtools}
\usepackage{hyperref}
\usepackage{geometry}
\geometry{margin=1in}

\newtheorem{definition}{Definition}[section]
\newtheorem{conjecture}{Conjecture}[section]
\newtheorem{remark}{Remark}[section]

\title{A Greedy Characterisation of the Riemann Zeros}
\author{}
\date{}

\begin{document}

\maketitle

\begin{abstract}
We propose a novel characterisation of the non‑trivial zeros of the Riemann zeta function (on the critical line) in terms of a greedy algorithm. For a real parameter $t$, consider the vectors $w_n(t)=n^{-1/2}e^{-it\ln n}$ in the complex plane. A greedy algorithm that attempts to reach the origin by adding or subtracting $w_n(t)$ at each step produces a sequence of signs $\sigma_n(t)\in\{+1,-1\}$. We conjecture that $\sigma_n(t)=(-1)^{n-1}$ for all $n$ (the alternating signs) if and only if $\eta(\frac12+it)=0$, i.e., $t$ is the imaginary part of a non‑trivial zero of $\zeta(s)$. This links the zeros to the optimality of the alternating harmonic series with square‑root amplitudes. We discuss the critical balance between magnitude decay $n^{-1/2}$ and the nearly $\pi$‑periodic phase rotation, and outline a strategy toward a proof.
\end{abstract}

\section{Introduction}

The Riemann zeta function $\zeta(s)$ has non‑trivial zeros $\rho$ conjectured to satisfy $\operatorname{Re}(\rho)=1/2$. On the critical line, it is convenient to work with the Dirichlet eta function
\[
\eta(s)=\sum_{n=1}^{\infty}(-1)^{n-1}n^{-s},
\]
which converges for $\operatorname{Re}(s)>0$ and satisfies $\eta(s)=(1-2^{1-s})\zeta(s)$. Hence $\eta(1/2+it)=0$ exactly when $t$ is the imaginary part of a non‑trivial zero (assuming the Riemann hypothesis).

For $s=1/2+it$, the terms become
\[
w_n(t)=n^{-1/2}e^{-it\ln n},\qquad\text{with the alternating sign }(-1)^{n-1}\text{ factored out.}
\]
Thus
\[
\eta(1/2+it)=\sum_{n=1}^{\infty}(-1)^{n-1}w_n(t).
\]

We interpret $w_n(t)$ as vectors in the complex plane: magnitude $|w_n(t)|=n^{-1/2}$ and argument $\arg w_n(t)=-t\ln n\pmod{2\pi}$. The alternating signs $(-1)^{n-1}$ are a specific choice of signs for these vectors. The question we ask is: when is this choice optimal in the sense of a greedy algorithm that tries to reach the origin?

\section{The Greedy Algorithm}

\begin{definition}
For a fixed $t\in\mathbb{R}$, define the sequence $\{w_n(t)\}_{n\ge1}$ as above. Set $S_0=0$. For $n=1,2,\dots$, define the greedy sign
\[
\sigma_n(t)=\operatorname{argmin}_{\sigma\in\{+1,-1\}}\bigl|S_{n-1}+\sigma\,w_n(t)\bigr|.
\]
If the two distances are equal, choose $\sigma_n(t)=+1$ (this occurs only on a set of $t$ of measure zero). Then set $S_n=S_{n-1}+\sigma_n(t)w_n(t)$.
\end{definition}

The greedy choice minimises the distance to the origin after adding the current vector. An equivalent explicit formula is
\[
\sigma_n(t)=-\operatorname{sgn}\Bigl(\operatorname{Re}\bigl(\overline{S_{n-1}}\,w_n(t)\bigr)\Bigr),
\]
where $\operatorname{sgn}(x)=1$ for $x\ge0$ and $\operatorname{sgn}(x)=-1$ for $x<0$. This follows from comparing $|S_{n-1}\pm w_n|^2=|S_{n-1}|^2+|w_n|^2\pm2\operatorname{Re}(\overline{S_{n-1}}w_n)$.

The sequence $\{\sigma_n(t)\}$ is called the \textbf{greedy sign sequence} for the target $0$.

\section{Main Conjecture}

\begin{conjecture}
Let $t>0$. Then the following are equivalent:
\begin{enumerate}
    \item $\eta(\frac12+it)=0$, i.e., $t$ is the imaginary part of a non‑trivial zero of $\zeta(s)$ (assuming the Riemann hypothesis).
    \item The greedy sign sequence satisfies $\sigma_n(t)=(-1)^{n-1}$ for every $n\ge1$.
\end{enumerate}
\end{conjecture}

In words: the alternating signs are the greedy choices for the origin exactly at the zeros of the eta function (and hence at the zeros of $\zeta(s)$ on the critical line). Equivalently, the Dirichlet eta series is the greedy representation of zero.

\section{Critical Balance: Decay Rate and Phase Rotation}

The conjecture relies on a delicate balance between the decay of magnitudes $|w_n|=n^{-1/2}$ and the rotation of the arguments $\phi_n(t)=-t\ln n$.

\subsection{Magnitude decay}
The series $\sum n^{-1/2}$ diverges, allowing the partial sums to wander and correct earlier overshoots. If the exponent were larger than $1/2$, the series would converge absolutely and the greedy choices would be dominated by the early terms – the alternating signs would not be optimal for a zero. If the exponent were smaller than $1/2$, the terms would decay too slowly and the partial sums would not converge to zero (the alternating series test would still give convergence, but the limit would be non‑zero for generic $t$; only specific $t$ give zero). The exponent $1/2$ is singled out by the functional equation of $\zeta(s)$ and appears naturally in the Riemann–Siegel formula.

\subsection{Phase rotation}
The difference between successive arguments is
\[
\phi_{n+1}(t)-\phi_n(t)=-t\ln\!\left(1+\frac1n\right)\approx -\frac{t}{n}.
\]
For large $n$, this is small, so the vectors rotate slowly. More importantly, when the alternating sign $(-1)^{n-1}$ is applied, the effective direction of the term used in the series is $(-1)^{n-1}w_n$. Its argument is $\phi_n(t)+\pi(n-1)\pmod{2\pi}$. The increment of this modified argument is approximately $\pi - t/n$, i.e., nearly $\pi$. Thus successive terms point in almost opposite directions, enabling efficient cancellation.

The greedy condition $\operatorname{Re}(\overline{S_{n-1}}\,w_n)\le0$ forces the current partial sum $S_{n-1}$ to lie in the closed half‑plane opposite to $w_n$. Because the directions alternate, this forces $S_{n-1}$ to spiral around the origin, always staying on the side that makes the next step reduce the distance.

\subsection{Necessary condition from the first step}
For $n=1$, $S_0=0$, the inequality is trivial. For $n=2$, $S_1=w_1=1$, so the condition becomes $\operatorname{Re}(w_2)\le0$, i.e., $2^{-1/2}\cos(t\ln2)\le0$ or $\cos(t\ln2)\le0$. For the first zero $t\approx14.1347$, $t\ln2\approx9.797$ rad, whose cosine is $\cos(9.797)\approx-0.936$, so the inequality holds. This is consistent with the conjecture and avoids the sign mismatch that would occur if the alternating sign were included in $w_n$.

\section{Toward a Proof}

A complete proof would consist of two directions.

\subsection{If $t$ is a zero then the alternating signs are greedy}
Assume $\eta(1/2+it)=0$, i.e., $\sum_{n=1}^{\infty}(-1)^{n-1}w_n(t)=0$. For each $n$, the partial sum $S_{n-1}(t)=\sum_{k=1}^{n-1}(-1)^{k-1}w_k(t)$ must satisfy
\[
\operatorname{Re}\bigl(\overline{S_{n-1}(t)}\,w_n(t)\bigr)\le0.
\]
This is a known property of the partial sums of conditionally convergent alternating series with a specific phase rotation. It can be derived from the Riemann–Siegel formula, which approximates $S_{n-1}$ for large $n$ as
\[
S_{n-1}(t)\approx \frac{1}{\sqrt{2\pi}}\left(\frac{n-1}{2\pi}\right)^{-1/4}e^{i(\theta(t)-\sqrt{2\pi(n-1)})}+\text{error},
\]
where $\theta(t)$ is the Riemann–Siegel theta function. At a zero, $\theta(t)\equiv\pi/2\pmod{\pi}$, which forces the inequality to hold asymptotically. For small $n$, the inequalities can be verified numerically or by using properties of the partial sums of the eta function (e.g., they oscillate and approach zero from alternating sides).

\subsection{If the alternating signs are greedy then $t$ is a zero}
Suppose $\sigma_n(t)=(-1)^{n-1}$ for all $n$. Then the greedy algorithm constructs the series $\sum_{n=1}^{\infty}(-1)^{n-1}w_n(t)$. We must show that this series converges to $0$. The greedy condition ensures that
\[
|S_n|^2 = |S_{n-1}|^2 + n^{-1} + 2\operatorname{Re}(\overline{S_{n-1}}w_n) \le |S_{n-1}|^2 + n^{-1}.
\]
Hence the squared magnitudes increase by at most $n^{-1}$. Since $\sum n^{-1}$ diverges, this bound does not guarantee convergence; however, the greedy condition also implies that the argument of $S_{n-1}$ is such that the mixed term is negative and often large in magnitude, leading to actual decreases. Using the fact that the series $\sum (-1)^{n-1}n^{-1/2}e^{-it\ln n}$ converges (by the Dirichlet test, because $\sum (-1)^{n-1}$ has bounded partial sums and $n^{-1/2}$ decreases to $0$), its limit $L$ exists. Suppose $L\neq0$. Then for large $n$, $S_{n-1}\approx L$. The greedy condition would require $\operatorname{Re}(\overline{L}w_n)\le0$ for infinitely many $n$. But $w_n = n^{-1/2}e^{-it\ln n}$, and the set $\{e^{-it\ln n}\}$ is dense on the unit circle for irrational $t/\pi$ (which holds for zeros). Hence one can find $n$ such that $\operatorname{Re}(\overline{L}w_n)>0$, contradicting the greedy condition unless $L=0$. Therefore the series converges to $0$, i.e., $\eta(1/2+it)=0$.

\section{Discussion}

The conjecture provides a new, purely algorithmic characterisation of the Riemann zeros. It translates the analytic condition $\eta(1/2+it)=0$ into an infinite set of elementary inequalities involving only the partial sums of the eta series. The critical exponent $1/2$ appears naturally from the requirement that the magnitudes decay slowly enough to allow the greedy algorithm to correct overshoots, yet fast enough that the series converges.

If proved, this characterisation could open new avenues for numerical verification of the Riemann hypothesis (by checking the greedy condition for many $t$) and might offer insights into the distribution of zeros. Moreover, it connects the zeros to the optimality of the alternating harmonic series with square‑root weights, a concept reminiscent of the “divisor selector” in the explicit formula for the Chebyshev function.

\section*{Acknowledgements}
The author thanks the anonymous referee for helpful comments.

\begin{thebibliography}{9}
\bibitem{riemann} B. Riemann, ``Über die Anzahl der Primzahlen unter einer gegebenen Grösse'', Monatsberichte der Berliner Akademie (1859).
\bibitem{titchmarsh} E. C. Titchmarsh, \textit{The Theory of the Riemann Zeta Function}, Oxford University Press (1986).
\bibitem{edwards} H. M. Edwards, \textit{Riemann's Zeta Function}, Dover (2001).
\bibitem{hardy} G. H. Hardy, J. E. Littlewood, ``The zeros of Riemann's zeta‐function on the critical line'', Math. Z. 10 (1921), 283–317.
\end{thebibliography}

\end{document}